%%%% CAPÍTULO 2 - REVISÃO DA LITERATURA (OU REVISÃO BIBLIOGRÁFICA, ESTADO DA ARTE, ESTADO DO CONHECIMENTO)
%%
%% O autor deve registrar seu conhecimento sobre a literatura básica do assunto, discutindo e comentando a informação já publicada.
%% A revisão deve ser apresentada, preferencialmente, em ordem cronológica e por blocos de assunto, procurando mostrar a evolução do tema.
%% Título e rótulo de capítulo (rótulos não devem conter caracteres especiais, acentuados ou cedilha)
\chapter{Referencial Teórico}\label{cap:referencialTeorico}

O aprendizado de máquina moderno envolve uma série de decisões interdependentes que impactam diretamente a eficácia dos modelos, desde o pré-processamento dos dados até a escolha e configuração dos algoritmos. A crescente complexidade desses processos impulsionou o surgimento do \textit{AutoML}, cujo objetivo é automatizar etapas críticas de um pipeline preditivo. No entanto, para que essa automação seja eficiente, é necessário compreender os fundamentos que orientam tais decisões.

Este capítulo apresenta os principais conceitos que embasam a proposta deste trabalho, abordando inicialmente os fundamentos do Aprendizado de Máquina e a estrutura dos sistemas de \textit{AutoML}. Em seguida, discute-se o meta-aprendizado como estratégia para reutilizar conhecimento adquirido em problemas anteriores, otimizando escolhas futuras. Também são exploradas as meta-características, essenciais para caracterizar conjuntos de dados e treinar metamodelos eficazes, bem como o aprendizado ativo, que visa otimizar o processo de rotulagem. Por fim, são apresentadas estratégias para seleção e estruturação de pipelines, os componentes de pré-processamento e classificação utilizados no experimento, as representações vetoriais de conjuntos de dados e as métricas de ranqueamento adotadas na avaliação.


\section{Aprendizado de Máquina e AutoML}
A aprendizagem de máquina envolve sistemas que aprendem automaticamente a partir de dados, sem necessidade de programação explícita do modelo. \citeonline{Domingos2012} destaca que, para não se perder na variedade de algoritmos disponíveis, deve-se encarar o aprendizado como a combinação de três componentes principais: representação, avaliação e otimização. Nesse contexto:

\begin{itemize}
  \item \textbf{Representação:} define a linguagem formal dos classificadores que o sistema pode aprender, isto é, o espaço de hipóteses permitido. Se um classificador não está nesse espaço, ele não pode ser obtido pelo aprendizado;
  \item \textbf{Avaliação:} função objetivo utilizada para distinguir classificadores bons de maus (por exemplo, acurácia, erro quadrático). Em geral, há uma função interna de otimização distinta da métrica externa desejada;
  \item \textbf{Otimização:} método de busca no espaço de hipóteses para encontrar o classificador de maior pontuação. Técnicas comuns incluem busca combinatória, algoritmos evolucionários, gradiente etc.
\end{itemize}

A automação desses componentes é o foco do \gls{automl}. Sistemas de AutoML têm se mostrado eficazes ao empregar \emph{metamodelos} — modelos treinados sobre o histórico de desempenho de algoritmos em execuções anteriores, e não diretamente sobre os dados do problema-alvo — de forma \emph{offline} (isto é, previamente, a partir de execuções já realizadas em outras bases) para guiar a seleção de algoritmos e parâmetros, abordagem detalhada na Seção \ref{sec:metaaprendizado}. Em particular, \citeonline{Schoenfeld2020} observa que sistemas AutoML obtêm melhor desempenho usando metamodelos offline para estimar o desempenho de modelos. \citeonline{Vanschoren2018} aponta também que o meta-aprendizado pode acelerar e aprimorar o projeto de pipelines de ML, substituindo abordagens manuais por soluções aprendidas de forma orientada por dados.

\section{Meta-aprendizado}\label{sec:metaaprendizado}
O meta-aprendizado, ou "aprender a aprender" \citeonline{Vilalta2002}, é a ciência de observar sistematicamente como diferentes abordagens de aprendizado de máquina performam em múltiplas tarefas, e então aprender com essa experiência (metadados) para resolver novas tarefas mais rapidamente, de modo a construir sistemas de aprendizado autoadaptativos que melhoram seu viés indutivo dinamicamente com a experiência acumulada \citeonline{Brazdil2008}. Em outras palavras, utiliza-se o histórico de desempenho de algoritmos em uma coleção de bases de dados para treinar um \emph{metamodelo} (modelo treinado sobre metadados de execuções anteriores, e não diretamente sobre os dados do problema-alvo) que pode prever o melhor algoritmo ou configuração para um novo problema. Adicionalmente, \citeonline{Schoenfeld2022} mostraram que, em vez de tratar modelos como caixas-pretas, pode-se explorar metadados granulares dos pipelines para melhorar a previsão de desempenho.

As abordagens de meta-aprendizado podem ser \emph{offline} ou \emph{online}. Em métodos \emph{offline}, diversos pipelines são executados previamente em coleções de conjunto de dados e os resultados formam o meta conjunto usado para treinar o metamodelo. Sistemas AutoML como o Auto-sklearn \citeonline{FeurerAutosklearn2015} e o SmartML \citeonline{Maher2019} ilustram esse caso, utilizando os resultados de múltiplas execuções em bases conhecidas para guiar a busca em novas bases. Em contrapartida, métodos \emph{online}, como o TPOT \cite{Olson2016TPOT}, atualizam o aprendizado à medida que geram pipelines no próprio problema de interesse.

Além disso, benchmarks específicos como o \textit{meta conjunto}, proposto por \citeonline{Triantafillou2020}, foram desenvolvidos para avaliar modelos de few-shot learning e meta-aprendizado em cenários mais realistas e desafiadores. O \textit{meta conjunto} é composto por um conjunto diversificado de conjuntos de dados que permite avaliar a capacidade dos meta-modelos de generalizar para novos domínios e tarefas com poucos exemplos. A proposta introduz variações no número de classes, tamanho de conjuntos de suporte e testes com desbalanceamento, promovendo uma avaliação mais robusta dos métodos de meta-aprendizado.

Em todos os casos, o objetivo é transferir conhecimento entre tarefas: utilizar metadados e meta características coletados em problemas anteriores para reduzir o tempo e esforço de modelagem em novos problemas.

\section{Meta-características e Caracterização de Conjuntos de Dados}
A caracterização de dados por meio de \emph{meta características} é essencial para o meta-aprendizado. Meta características são atributos que descrevem propriedades estatísticas e estruturais de um conjunto de dados, e servem de entrada para os metamodelos. Por exemplo, \citeonline{Schoenfeld2022} extraíram 72 meta características de cada base de dados, incluindo categorias simples, estatísticas, informações teóricas, \emph{landmarking} e baseadas em modelo. Esses grupos podem ser definidos genericamente da seguinte forma:

\begin{itemize}
  \item \textbf{Simples:} propriedades elementares do conjunto, como número de instâncias, número de atributos, proporção de classes;
  \item \textbf{Estatísticas:} medidas como médias, variâncias e correlações dos atributos;
  \item \textbf{Informacionais:} métricas de teoria da informação, como entropia das classes e entropia dos atributos;
  \item \textbf{Landmarking:} desempenho de algoritmos simples e rápidos (por exemplo, taxa de erro de 1-vizinho mais próximo, árvore de decisão rasa) no conjunto de dados;
  \item \textbf{Baseadas em modelo:} características derivadas de modelos treinados, como profundidade média de árvores de decisão, coeficientes de regressão, etc.
\end{itemize}

Essas meta características capturam informações relevantes que permitem medir similaridades entre conjuntos de dados e orientar o aprendizado \citeonline{Feurer2015}. Muitas pesquisas utilizam repositórios de dados públicos; por exemplo, o conjunto de 194 bases usado por \citeonline{Schoenfeld2022} provém do programa \gls{d3m} (\textit{Data-Driven Discovery of Models}), iniciativa da \gls{darpa} voltada à automação da ciência de dados \citeonline{Lippmann2016}, e inclui conjuntos de dados originados de repositórios como o OpenML \citeonline{Vanschoren2014} e o \gls{uci} Machine Learning Repository \cite{Dua2019}. Cada instância do meta conjunto consiste no vetor de metacaracterísticas do conjunto de dados e na estrutura do pipeline avaliado, associado ao respectivo desempenho (por exemplo, medida F1) obtido naquela base. Esse meta conjunto (meta características + descrição de pipeline + desempenho) é então utilizado para treinar metamodelos capazes de prever ou ranquear pipelines em novos conjuntos de dados.

\section{Aprendizado Ativo}

O aprendizado ativo é uma abordagem em que o algoritmo seleciona interativamente quais instâncias devem ser rotuladas para melhorar seu desempenho com poucos dados rotulados \citeonline{Cohn1994}. \citeonline{Settles2009} descreve que a ideia central do aprendizado ativo é permitir que o sistema escolha os exemplos a serem rotulados, de forma que se atinja maior precisão com um número reduzido de etiquetas. Em outras palavras, um modelo com capacidade de consulta pode alcançar desempenho elevado utilizando apenas os exemplos mais informativos. Esse paradigma é motivado por cenários em que há abundância de dados não rotulados e o custo de rotular é alto.

Para implementar o aprendizado ativo, definem-se estratégias de consulta que estimam a utilidade de rotular cada instância disponível. Abordagens comuns incluem amostragem por incerteza (rotular as instâncias nas quais o modelo está menos confiante), seleção por comitê (vários modelos votam nas instâncias mais controversas) e redução de erro esperado.

No contexto de meta-aprendizado para AutoML, o aprendizado ativo é relevante como extensão futura: ao invés de avaliar todos os 60 pipelines em cada novo conjunto de dados, um sistema poderia selecionar ativamente quais pipelines avaliar primeiro com base na incerteza do metamodelo, reduzindo o custo computacional da construção do meta-conjunto. \citeonline{douce2024bridging} e \citeonline{Sener2018} exemplificam essa abordagem ao utilizar \textit{embeddings} para guiar a seleção de exemplos informativos no espaço de características — princípio análogo ao uso de \textit{embeddings} de conjuntos de dados para identificar quais avaliações de pipeline são mais informativas para o sistema de recomendação.

\section{Seleção e estrutura de \textit{pipelines}}
Em sistemas AutoML, um \emph{pipeline} de aprendizado de máquina consiste na composição de múltiplos passos de pré-processamento e modelagem que transformam os dados até a predição final. Em geral, pipelines são representados como grafos acíclicos direcionados (\glsadd{dag}DAGs), podendo incluir estruturas paralelas como comitês. Essa flexibilidade permite representar desde um pipeline simples (uma sequência linear de transformação e classificador) até configurações complexas (ex.: três sub pipelines com combinação por votação majoritária).  

Entretanto, o número de pipelines possíveis cresce rapidamente com o número de algoritmos e passos disponíveis. Como observado por \citeonline{Schoenfeld2022}, as diversas maneiras de combinar algoritmos em pipelines criam um espaço de busca \emph{não escalonável} de possíveis fluxos a serem considerados. Diante disso, é útil empregar estratégias de pré-seleção ou ranqueamento de pipelines candidatos para priorizar a execução daqueles com maior potencial (especialmente quando recursos computacionais são limitados). Por exemplo, abordagens tradicionais de AutoML, como Auto-sklearn \cite{FeurerAutosklearn2015} e SmartML \cite{Maher2019}, tratam o pipeline como uma caixa-preta e utilizam apenas o desempenho observado e a similaridade de meta características do conjunto de dados para ranquear modelos. Outras soluções simplificam a estrutura do pipeline, por exemplo considerando apenas sequências lineares e usando uma rede \gls{lstm} para estimar valor\footnote{LSTM: rede neural recorrente que aprende dependências de longo prazo em sequências por meio de um mecanismo de portas que regula o fluxo de informação em uma célula de memória interna \citeonline{Hochreiter1997}.}.

Em contraste, trabalhos mais recentes exploram explicitamente a estrutura do pipeline. \citeonline{Schoenfeld2020} propõe um metamodelo baseado em rede neural dinamicamente montada, capaz de modelar DAGs arbitrários do pipeline. \citeonline{Schoenfeld2022} validam empiricamente essa abordagem criando um grande meta conjunto (cerca de 400.000 instâncias de aproximadamente 4.000 pipelines avaliados em 194 bases) e mostram que metamodelos que consideram metadados granulares de cada pipeline superam aqueles que usam representações simplificadas. Concretamente, eles demonstraram que metamodelos que utilizam toda a estrutura do pipeline fornecem estimativas de desempenho mais precisas, e que ranquear pipelines baseado na similaridade das meta características do conjunto de dados resulta em recomendações superiores às baseadas em estimativas de desempenho absolutas. 

Em resumo, a seleção de pipelines em AutoML combina conhecimento do domínio (via meta-características de dados) e exploração da estrutura dos fluxos de processamento. Esse campo de pesquisa continua em desenvolvimento, mas já existe consenso de que considerar explicitamente a granularidade e a estrutura dos pipelines é crucial para sistemas AutoML eficientes e escaláveis.

\section{Componentes do Arsenal Experimental}
\label{sec:componentesArsenal}

Como a metodologia deste trabalho avalia combinações de escalonamento, seleção de atributos e classificação, é necessário explicitar o papel dos operadores que compõem o arsenal experimental. Todos os componentes foram implementados por meio da biblioteca \textit{scikit-learn}, referência consolidada para experimentação reprodutível em aprendizado de máquina com Python \cite{Pedregosa2011}.

No estágio de escalonamento, foram considerados três transformadores e uma alternativa sem transformação. O \texttt{StandardScaler} centraliza cada atributo pela média e o reescala pelo desvio padrão, tornando os atributos comparáveis em termos de variância. O \texttt{MinMaxScaler} reescala os valores para um intervalo fixo, usualmente $[0,1]$. O \texttt{RobustScaler} utiliza mediana e intervalo interquartil, sendo menos sensível a valores extremos. Esses operadores são relevantes porque algoritmos baseados em distância, margem ou otimização numérica, como \gls{knn}, \gls{svm} e \gls{mlp}, podem ser fortemente afetados por diferenças de escala entre atributos.

No estágio de seleção ou redução de atributos, foram utilizados \gls{pca}, \texttt{SelectKBest} e ausência de seleção. A \gls{pca} projeta os dados em componentes ortogonais que concentram a maior variância possível, sendo uma técnica clássica de redução de dimensionalidade \cite{Jolliffe2002}. O \texttt{SelectKBest}, neste trabalho, utiliza o critério \texttt{f\_classif}, baseado no teste F da \gls{anova}, para manter os atributos mais associados ao alvo; trata-se de uma estratégia de filtragem supervisionada de atributos \cite{Guyon2003}. A alternativa \texttt{None} preserva todos os atributos originais, permitindo avaliar quando a redução dimensional prejudica ou favorece o desempenho.

No estágio de classificação, foram avaliados cinco algoritmos. O \gls{rf} combina múltiplas árvores de decisão treinadas sobre amostras e subconjuntos de atributos, reduzindo variância e aumentando robustez em dados tabulares \cite{Breiman2001}. A \gls{svm} busca hiperplanos de máxima margem e, com kernel RBF, modela fronteiras não lineares \cite{Cortes1995}. O \gls{knn} é um método baseado em instâncias que classifica uma nova amostra pela vizinhança mais próxima no espaço de atributos \cite{Cover1967}. A \gls{mlp} é uma rede neural feedforward capaz de modelar relações não lineares por camadas ocultas \cite{Bishop2006}. Por fim, a \gls{dt} particiona recursivamente o espaço de atributos por regras de decisão, oferecendo interpretabilidade e baixo custo computacional \cite{Quinlan1986}.

%% ────────────────────────────────────────────────────────────────────────────
\section{Embeddings de Conjuntos de Dados e Métricas de Ranking}
\label{sec:embeddingsRanking}

Esta seção apresenta os dois fundamentos diretamente empregados na proposta deste trabalho: as representações vetoriais de conjuntos de dados conhecidas como \textit{embeddings}, utilizadas para caracterizar os conjuntos de dados do meta-conjunto, e as métricas empregadas para avaliar a qualidade de um ranqueamento de pipelines.

\subsection{Embeddings de Conjuntos de Dados}

Uma limitação das meta-características tradicionais é que cada conjunto de dados é representado por um único vetor de estatísticas globais — médias, variâncias, entropia das classes —, o que descarta a variabilidade interna da distribuição dos atributos. Para superar essa limitação, abordagens mais recentes propõem \textit{embeddings} de conjuntos de dados: vetores densos de dimensão fixa que capturam propriedades representativas do conjunto de dados de forma invariante ao número de atributos e instâncias.

\citeonline{Jomaa2021} propõem o Dataset2Vec, um \textit{embedding} neural baseado em arquitetura DeepSet que representa o conjunto de dados como um conjunto hierárquico de pares preditores-alvo. O modelo é treinado com uma tarefa auxiliar de similaridade entre conjuntos de dados e produz vetores de 64 dimensões cujas proximidades no espaço euclidiano correlacionam-se com a similaridade de desempenho de algoritmos. Os autores identificam quatro critérios desejáveis para embeddings de conjuntos de dados, a saber: (D1) independência de esquema — aplicável a qualquer conjunto sem restrição de dimensionalidade; (D2) expressividade — capacidade de distinguir conjuntos com características de aprendizado distintas; (D3) escalabilidade — calculável sem treinamento para novos conjuntos; (D4) correlação com meta-alvos — os vetores devem predizer desempenho de algoritmos.

Uma alternativa sem necessidade de treinamento consiste nos \textit{embeddings} estatísticos distribucionais: ao invés de uma única estatística global por atributo, calcula-se a distribuição desta estatística sobre todos os atributos do conjunto de dados via percentis, produzindo um vetor de dimensão fixa que satisfaz os critérios D1 e D3 e se aproxima de D2 em tarefas tabulares. Esta abordagem é análoga ao MF2\footnote{O MF2 é um conjunto de 46 meta-características estáticas — incluindo medidas de \textit{landmarking}, componentes principais (PCA) e proporção de valores ausentes — originalmente proposto por \citeonline{Reif2012} e empregado por \citeonline{Feurer2015} para estimar a similaridade entre conjuntos de dados na inicialização de otimização bayesiana de hiperparâmetros.} de \citeonline{Feurer2015}, porém com maior riqueza representacional por capturar distribuições e não apenas médias globais.

\subsection{Métricas de Avaliação de Ranking}

Na tarefa de ranqueamento de pipelines, a qualidade das predições é medida por métricas que comparam a ordenação predita com a ordenação real. As principais métricas utilizadas na literatura de meta-aprendizado são:

\textbf{Coeficiente de Correlação de Spearman ($\rho$):} mede a correlação monotônica entre os postos atribuídos aos scores preditos e reais. Valores próximos de 1 indicam que o ranking predito preserva a ordem correta dos pipelines; valores próximos de zero indicam ausência de correlação \citeonline{Spearman1904}.

\textbf{Kendall Tau ($\tau$):} mede a fração líquida de pares de pipelines corretamente ordenados. Dado um par $(i, j)$, o par é concordante se o ranking predito e o real concordam sobre qual pipeline é superior; é discordante caso contrário. O valor de $\tau$ é a diferença entre pares concordantes e discordantes, normalizada pelo total de pares \citeonline{Kendall1938}.

\textbf{Kendall Tau Ponderado ($\tau_w$):} extensão do Kendall Tau que atribui pesos maiores às inversões de ordem que ocorrem no topo do ranking. Implementado em \texttt{scipy.stats.weightedtau}, o peso de cada par de posições $(i,j)$ é inversamente proporcional ao produto de seus postos no ranking real, de modo que erros envolvendo os melhores pipelines são penalizados mais severamente do que erros nas posições inferiores.

\textbf{\glsadd{ndcg}NDCG} (\textit{Normalized Discounted Cumulative Gain}): métrica proveniente da área de recuperação de informação, que pondera a relevância dos itens recuperados pelo logaritmo da posição na lista. Em meta-aprendizado, o NDCG penaliza explicitamente pipelines de alto desempenho recomendados em posições distantes do topo.

A adoção conjunta dessas quatro métricas permite avaliar o sistema tanto em termos de correlação global (Spearman e Kendall Tau) quanto de precisão nas posições de maior impacto prático. Nesse segundo grupo, Kendall Tau ponderado e NDCG são métricas complementares: a primeira penaliza inversões pareadas no topo do ranking, enquanto a segunda mede a relevância acumulada das primeiras posições.

%% ────────────────────────────────────────────────────────────────────────────
\section{Trabalhos Relacionados}
\label{sec:trabalhosRelacionados}

Esta seção apresenta os principais trabalhos relacionados à proposta desta pesquisa, organizados em três eixos temáticos: (i) abordagens para extração de meta-características e representação de conjuntos de dados; (ii) sistemas de recomendação de algoritmos e pipelines de aprendizado de máquina via meta-aprendizado; e (iii) estratégias de avaliação de complexidade de fronteiras de decisão e aprendizado ativo. Ao final, é apresentada uma síntese comparativa que posiciona a presente proposta em relação ao estado da arte.

\subsection{Representação de Conjuntos de Dados e Meta-características}
\label{sec:relMetaFeatures}

A qualidade dos metamodelos em sistemas de meta-aprendizado depende diretamente da capacidade de representar conjuntos de dados de forma compacta e informativa. Historicamente, duas correntes principais se desenvolveram: a engenharia de meta-características estáticas e a aprendizagem de meta-características por redes neurais.

\citeonline{Vanschoren2014} introduzem o OpenML, uma plataforma colaborativa que centraliza conjuntos de dados, experimentos e resultados de algoritmos em um repositório público e estruturado. Além de prover acesso unificado a centenas de conjuntos de dados anotados, o OpenML define uma infraestrutura para coleta sistemática de meta-dados — acurácias, tempos de execução, hiperparâmetros utilizados — que viabiliza a construção de grandes meta-conjuntos. Esta plataforma constitui a principal fonte de dados do presente trabalho, tendo fornecido os 201 conjuntos utilizados nos experimentos.

\citeonline{Schoenfeld2022} investigam a construção de meta-conjuntos de grande escala para otimização de pipelines, criando uma base com aproximadamente 400 mil instâncias derivadas de 4 mil pipelines avaliados em 194 conjuntos do programa D3M. O trabalho demonstra empiricamente que metamodelos que consideram metadados granulares de toda a estrutura do pipeline superam abordagens que utilizam apenas representações simplificadas. Os autores mostram, ainda, que o ranqueamento baseado em similaridade de meta-características fornece recomendações superiores às estimativas de desempenho absoluto, convergindo diretamente com a abordagem adotada neste TCC.

\citeonline{Jomaa2021} propõem o \textit{Dataset2Vec}, um extrator de meta-features baseado em uma arquitetura DeepSet que representa conjuntos de dados como conjuntos hierárquicos de pares preditores-alvo. O modelo é treinado com uma tarefa auxiliar de aprendizagem de similaridade entre conjuntos de dados (\textit{dataset similarity learning}), que determina se dois subconjuntos amostrados provêm do mesmo conjunto. Os resultados experimentais com 120 conjuntos UCI para a tarefa de otimização de hiperparâmetros demonstram que as meta-features aprendidas pelo Dataset2Vec superam meta-features estáticas tradicionais (MF1 de \citeonline{Wistuba2016} com 22 dimensões; MF2 de \citeonline{Feurer2015} com 46 dimensões) e o \textit{Neural Statistician} (\citeonline{Edwards2017}) em termos de correlação de Spearman com os alvos de desempenho.

O presente trabalho diferencia-se do Dataset2Vec em dois aspectos fundamentais: (1) ao invés de embeddings neurais, utiliza-se um embedding distribucional de 64 dimensões, produzido sem necessidade de treinamento prévio, com propriedades análogas ao MF2 mas com maior riqueza por capturar a \textit{distribuição} das meta-features por atributo (não apenas médias globais); (2) a tarefa meta-nível é o \textit{ranqueamento de pipelines} (não apenas otimização de hiperparâmetros), sobre um meta-conjunto de 201 conjuntos de dados. Esta escolha de embedding estático garante a propriedade de independência de esquema (\textbf{D1} na taxonomia do Dataset2Vec), escalabilidade sem treinamento (\textbf{D3}) e reprodutibilidade, ao custo de expressividade potencialmente menor (\textbf{D2}) em relação a abordagens neurais.

\subsection{Sistemas de Recomendação de Algoritmos e Pipelines}
\label{sec:relRecomendacao}

\citeonline{Schoenfeld2020} propõem um metamodelo baseado em rede neural dinamicamente configurável, capaz de modelar DAGs arbitrários de pipelines. O sistema representa cada operador do pipeline com uma subrrede específica e combina as representações por uma função de pooling global. Esta abordagem permite estimar o desempenho de configurações de pipeline inéditas sem avaliação direta, tornando a exploração do espaço de busca mais eficiente. A arquitetura supera metamodelos com representações simplificadas de pipeline, reforçando a importância de considerar a estrutura interna dos fluxos de processamento.

\citeonline{ali2020mega} apresentam o sistema \gls{mega} (\textit{Meta-learning Ensemble for Genetic Algorithm}), um método de combinação de classificadores que emprega meta-aprendizado e algoritmo genético para selecionar a melhor combinação de classificadores para um conjunto de dados desconhecido. O sistema extrai meta-características dos conjuntos, usa um algoritmo genético para encontrar combinações eficazes e, na fase de inferência, aplica um k-NN ponderado para recomendar combinações ideais. O uso do k-NN como metamodelo de inferência é compatível com a abordagem adotada neste trabalho, onde também se emprega k-NN para predição de rankings de pipelines.

\citeonline{DBLP:journals/corr/abs-2106-03954} investigam a recomendação de algoritmos de seleção de atributos por meta-aprendizado, modelando a tarefa como um problema supervisionado. O trabalho emprega k-NN, SVM, árvores de decisão e florestas aleatórias como meta-aprendizes, com meta-características extraídas de 28 meta-conjuntos do repositório ASlib. Os autores demonstram que técnicas de redução de dimensionalidade (PCA, seleção por variância, Chi-quadrado, ANOVA-F) permitem reduzir em mais de 80\% o número de meta-características sem degradar o desempenho preditivo, evidenciando a robustez de abordagens baseadas em similaridade para recomendação.

\citeonline{wang2021ensemblelearningbasedclassification} propõem um método multi-rótulo de recomendação de algoritmos de classificação baseado em aprendizado por comitê. O modelo mapeia meta-características de 1090 problemas de classificação para o conjunto de algoritmos adequados, combinando múltiplos tipos de meta-características em modelos-base com integração ponderada. Os experimentos demonstram superioridade em métricas de ranking (perda de ordenação e \textit{Precision@m}) sobre modelos individuais, reforçando o valor de abordagens de comitê para recomendação.

\citeonline{khan} descrevem o AutoFE-Sel, sistema de meta-aprendizado para recomendação de algoritmos de seleção de subconjunto de atributos. A abordagem utiliza a métrica EARR (\textit{Adjusted Ratio of Ratios}), que combina acurácia, tempo de execução e número de atributos selecionados para comparar algoritmos. Experimentos com 125 conjuntos de dados e 12 algoritmos candidatos demonstram que o AutoFE-Sel supera abordagens tradicionais (MCFA, IBL-Ranking, Fuzzy-Sim) em acurácia, \textit{hit-ratio} e \textit{F-measure}, atingindo acerto superior a 92\%.

\subsection{Avaliação de Fronteiras de Decisão e Aprendizado Ativo}
\label{sec:relFronteirasAtivo}

\citeonline{varshney2015persistent} utilizam topologia persistente para caracterizar a complexidade das fronteiras de decisão em classificadores supervisionados, construindo complexos simpliciais a partir de amostras com rótulos distintos. Os números de Betti obtidos pela homologia persistente correlacionam-se com o desempenho preditivo dos modelos de SVM, auxiliando na seleção de hiperparâmetros. Embora a abordagem topológica seja distinta da proposta neste TCC, o princípio de usar características intrínsecas dos dados para orientar a seleção de modelos é compartilhado.

\citeonline{ramamurthy2018topological} expandem a análise topológica propondo os complexos \textit{labeled Čech} e \textit{labeled Vietoris-Rips} para estimar a complexidade topológica da fronteira de decisão (\gls{dbtc}). Experimentos em \gls{mnist}, \gls{fashionmnist} e \gls{cifar10} mostram que a correspondência entre a DBTC do modelo e a do conjunto de dados favorece a seleção de modelos pré-treinados, evidenciando que propriedades estruturais dos dados podem guiar a escolha de algoritmos.

\citeonline{DBLP:journals/corr/abs-2004-11149} apresentam uma revisão abrangente sobre meta-aprendizado, sistematizando métodos baseados em otimização, modelo, métrica e inferência Bayesiana. O trabalho discute como o aprendizado meta-episódico permite generalização eficiente para novas tarefas com dados rotulados limitados, especialmente em cenários de \textit{few-shot learning} (aprendizado a partir de pouquíssimos exemplos rotulados por classe), formalizando a notação de meta-tarefas $\tau_i \sim \mathcal{P}(\mathcal{T})$ com conjuntos de suporte $\mathcal{D}^s_i$ e consulta $\mathcal{D}^q_i$.

\citeonline{douce2024bridging} propõem o método \gls{tcm}, que combina amostragem baseada em diversidade e incerteza para aprendizado ativo. O TCM realiza uma transição gradual de \textit{TypiClust} — que seleciona exemplos representativos por clustering em embeddings auto-supervisionados — para \textit{Margin Sampling}, focado na incerteza do modelo. Os autores demonstram superioridade sobre \gls{badge}, Coreset e ProbCover em \gls{cifar10}, \gls{cifar100} e \gls{isic2019}, evidenciando o valor de combinar diversidade e incerteza na seleção de instâncias informativas. A utilização de embeddings para identificar exemplos representativos é análoga ao uso de embeddings de conjuntos de dados neste trabalho para identificar conjuntos similares no espaço de meta-aprendizado.

\subsection{Síntese Comparativa}
\label{sec:sintese-comparativa}

A \autoref{tab:comparativo} sintetiza os principais trabalhos relacionados em relação a quatro dimensões: tipo de representação dos dados, tarefa de meta-nível, tamanho do meta-conjunto e método de avaliação utilizado.

\begin{table}[H]
\centering
\caption{Comparativo entre trabalhos relacionados e a proposta deste TCC}
\label{tab:comparativo}
\resizebox{\textwidth}{!}{%
\begin{tabular}{@{}lllll@{}}
\toprule
\textbf{Trabalho} & \textbf{Representação} & \textbf{Tarefa Meta-nível} & \textbf{Meta-conjunto} & \textbf{Avaliação} \\
\midrule
Schoenfeld et al. (2022)  & Estrutura do pipeline (DAG)    & Estimativa de desempenho       & $\sim$400 k instâncias & Correlação de Spearman \\
Dataset2Vec (2021)        & Embedding neural (DeepSet)     & Otimização de hiperparâmetros  & 120 conjuntos UCI      & Spearman (ranqueamento)  \\
MEGA (2020)               & Meta-características estáticas & Recomendação de comitês        & Múltiplos conjuntos    & Acurácia                 \\
AutoFE-Sel (2021)         & Meta-características estáticas & Seleção de atributos           & 125 conjuntos          & EARR, hit-ratio          \\
Wang et al. (2021)        & Meta-características multi-tipo & Recomendação multi-rótulo      & 1090 conjuntos         & Ranking loss, Prec@m      \\
\textbf{Este TCC}         & \textbf{embedding distribucional 64D} & \textbf{Ranqueamento de pipelines} & \textbf{201 conjuntos / 60 pipelines} & \textbf{Spearman, Kendall ($\tau$, $\tau_w$), NDCG} \\
\bottomrule
\end{tabular}
}
\fonte{}
\end{table}

A análise comparativa revela que a presente proposta distingue-se dos trabalhos anteriores principalmente por três aspectos: (1) foco explícito no \textit{ranqueamento completo de pipelines} (não apenas seleção do melhor ou estimativa pontual); (2) adoção conjunta de \textbf{Kendall Tau ponderado} (\textit{weightedtau}) e NDCG como métricas complementares de topo, alinhando a avaliação às necessidades reais de AutoML, onde os primeiros pipelines recomendados têm maior impacto prático; e (3) análise estatística da correlação entre dimensões do embedding e qualidade de ranqueamento por conjunto de dados, investigando \textit{quando} e \textit{por que} o meta-aprendizado funciona melhor.

A escolha de uma representação estatística distribucional, referida ao longo do trabalho como embedding distribucional por produzir um vetor denso de dimensão fixa, foi motivada pela compatibilidade com Python 3.13+, onde TensorFlow não possui distribuição, e pelo objetivo de estabelecer um baseline reprodutível e publicável. Diferentemente de embeddings neurais como o Dataset2Vec, essa representação não é aprendida por otimização de uma função objetivo; ela se aproxima do MF2 de Feurer et al. (2015) em termos de engenharia de meta-características, mas captura distribuições de meta-features por atributo em vez de apenas estatísticas globais do conjunto de dados.
